\chapter{The 4D Architecture: Event-Sourced Knowledge Graphs}
\label{ch:architecture_4d}

\section{Introduction to 4D Knowledge Representation}

Traditional knowledge graphs operate in a three-dimensional space of subjects, predicates, and objects. While powerful for representing static relationships, they often fail to capture the temporal evolution of facts without complex reification or versioning schemes. This chapter introduces the 4D Architecture: a paradigm shift that treats time and causality as first-class dimensions within the knowledge graph, implemented through event-sourced semantics and delivered via modular \textit{Ontology Packs}.

\section{Innovation 1: Event-Sourced Knowledge Graphs}

The core of the 4D architecture is the transformation of the RDF store from a mutable snapshot into an immutable, append-only event log.

\subsection{Immutable Append-Only Semantics}
Unlike traditional RDF stores (e.g., Apache Jena, GraphDB) that maintain the current state, the 4D Event-Sourced Knowledge Graph persists every state transition as a discrete event. Each event contains RDF deltas in canonical N-Quads format, ensuring that the entire history of the graph is cryptographically verifiable and immutable.

\begin{algorithm}
\caption{Append Event with RDF Deltas}
\begin{algorithmic}
  \State $eventId \gets$ UUID()
  \State $timestamp \gets$ now() \Comment{nanosecond precision}
  \State $vectorClock \gets$ vectorClock.increment()
  \State $deltaQuads \gets$ serialize(deltas) \Comment{N-Quads canonical format}
  \State quad $\gets$ (eventId, rdf:type, event:Event, EventLog)
  \State add(quad) \Comment{to EventLog named graph}
  \State \textbf{for each} $delta \in deltas$ \textbf{do}
    \State process(delta) \Comment{apply to Universe graph}
    \State add((eventId, event:hasDelta, $\cdot$, EventLog))
  \State \textbf{end for}
  \State \textbf{return} receipt $\gets$ \{eventId, timestamp, vectorClock, eventCount\}
\end{algorithmic}
\end{algorithm}

\subsection{Content-Addressable Integrity}
By utilizing N-Quads canonical ordering, we enable content-addressing for individual events. This allows for cryptographic integrity checks (e.g., BLAKE3 or SHA-256) across the event stream, providing a foundation for audit compliance and trust in distributed environments.

\section{Innovation 2: Vector Clock Causality Tracking}

In a distributed knowledge system, wall-clock time is insufficient for determining the order of events. The 4D architecture incorporates \textit{Vector Clocks} to track causal relationships across heterogeneous nodes.

\subsection{Causal Ordering across Distributed Nodes}
Every event in the 4D log carries a vector clock encoding its causal history:
\begin{equation}
VC = \{node_1: t_1, node_2: t_2, \ldots, node_n: t_n\}
\end{equation}

This enables the system to distinguish between sequential events and concurrent updates that may require conflict resolution. In the context of a swarm-native system, vector clocks provide the necessary synchronization primitive for maintaining a globally consistent (or eventually consistent) state without a central authority.

\section{Innovation 3: Deterministic 4D Time-Travel Reconstruction}

The addition of the fourth dimension (time) enables \textit{Time-Travel Reconstruction}: the ability to reconstruct the exact state of the knowledge graph at any arbitrary timestamp $t$ in the past.

\subsection{The Reconstruction Algorithm}
Reconstruction is performed by locating the nearest historical snapshot $s \le t$ and replaying the causal event log up to the target timestamp.

\begin{algorithm}
\caption{Deterministic State Reconstruction}
\begin{algorithmic}
  \State $snapshot \gets$ findLatestSnapshot($targetTime$)
  \State $reconstructed \gets$ loadSnapshot($snapshot$)
  \State $events \gets$ query(EventLog, $\{$ timestamp $\leq targetTime, timestamp > snapshot.time \}$)
  \State sortByVectorClock($events$) \Comment{Causal ordering}
  \State \ForEach{$event \in events$}
    \State applyDeltas($reconstructed, event.deltas$)
  \EndFor
  \State \Return $reconstructed$
\end{algorithmic}
\end{algorithm}


\subsection{Performance and Determinism}
By leveraging cached snapshot pointers in the metadata layer and streaming event replay, the 4D architecture achieves sub-5-second reconstruction for typical enterprise datasets. Because the reconstruction is pure and functional, it is guaranteed to be deterministic, yielding the same cryptographic hash for the same timestamp across all nodes.

\section{Ontology Packs: The Modular Delivery Mechanism}

To manage the complexity of 4D datasets and their associated logic, we introduce \textit{Ontology Packs} as the standard delivery mechanism.

\subsection{Defining Ontology Packs}
An Ontology Pack is a specialized bundle containing:
\begin{itemize}
    \item \textbf{RDF/OWL Ontologies}: The semantic definitions of the domain.
    \item \textbf{Validation Hooks}: Sub-microsecond policy functions for data integrity.
    \item \textbf{Code Templates}: Configurations for generating type-safe bindings in Rust, TypeScript, or Erlang.
\end{itemize}

\subsection{Integration with 4D Architecture}
Ontology Packs provide the metadata configuration required to interpret the 4D event stream. They define the namespaces, prefixes, and structural constraints (SHACL) that the reconstruction engine uses to validate state transitions. By modularizing these definitions, organizations can version and deploy temporal schemas as discrete, package-managed units.

\begin{lstlisting}[language=toml, caption=Example Ontology Pack Configuration]
[ontology]
default_namespace = "https://fusion.io/v1/"

[[ontology.ontologies]]
id = "core-4d"
name = "Core 4D Vocabulary"
version = "1.2.0"
file_path = "ontologies/4d.ttl"
format = "turtle"

[[ontology.targets]]
language = "rust"
features = ["oxigraph", "serde"]
template_path = "templates/rust/model.tera"
\end{lstlisting}

\section{Conclusion}

The 4D architecture moves beyond static data models by integrating event sourcing, vector clocks, and deterministic reconstruction into a unified semantic framework. Delivered through Ontology Packs, this architecture provides the robustness required for high-stakes enterprise applications while maintaining the flexibility of the Open World Assumption. In the following chapters, we will explore how this 4D foundation enables autonomic swarms to operate with information-theoretic correctness.

